% --- LaTeX Essay Template - S. Venkatraman ---

% --- Set document class and font size ---

\documentclass[letterpaper, 11pt]{article}

% --- Package imports ---

% lipsum is just for generating placeholder text and can be removed
\usepackage{hyperref, lipsum} 

% --- Fonts (Set to Palatino or Times New Roman if desired) ---

% \usepackage{newpxtext, newpxmath}
% \usepackage{times}

% --- Page layout settiings ---

% Set page margins
\usepackage[left=1.2in, right=1.2in, top=.9in, bottom=.9in]{geometry}

% Anchor footnotes to the bottom of the page
\usepackage[bottom]{footmisc}
\usepackage{setspace}
\setstretch{1.15}
% Set line spacing
\renewcommand{\baselinestretch}{1.2}

% Use this for a one-off '*' footnote (e.g., a URL in a heading)
\newcommand{\starfootnote}[1]{%
\begingroup
\renewcommand{\thefootnote}{\fnsymbol{footnote}}%
\footnote{#1}%
\endgroup
\addtocounter{footnote}{-1}%
}

% Set spacing between paragraphs
\setlength{\parskip}{2mm}

% --- Page formatting settings ---



% --- Essay title and author ---
\title{\vspace{-1.5cm} Summary of \textit{Five ways to ensure models serve society}\footnote{\url{https://www.nature.com/articles/d41586-020-01812-9}}}

\date{\today}

% --- Main text ---
% --- Document starts here ---

\begin{document}
% Set link colors for labeled items and URLs (blue) 
\hypersetup{colorlinks=true, linkcolor=blue, urlcolor=blue}
\maketitle
The article argues that models should serve society with humility, transparency, and awareness of uncertainty rather than political convenience.
It tells us five things to mind: assumptions, hubris, framing, consequences, and unknowns.
Mathematical models could be powerful tools for understanding complex systems and informing policy decisions, but they are not infallible.
Instead of treating models as objective truth, we should examine uncertainty, clarify value judgments, and be transparent about the limitations of our models.
Numbers could be manipulated to support any argument, and could be used to justify political agendas.
We should be cautious about the assumptions we make when building models, and be aware of the potential for bias and overconfidence.

\end{document}
